\documentclass[conference]{IEEEtran}

% 中文支持
\usepackage[UTF8]{ctex}

% 数学公式
\usepackage{amsmath,amsfonts,amssymb}

% 图表
\usepackage{graphicx}
\usepackage{booktabs}
\usepackage{float}
\usepackage{multirow}

% 超链接
\usepackage{hyperref}
\usepackage{xcolor}
\hypersetup{
    colorlinks=true,
    linkcolor=blue,
    citecolor=blue,
    urlcolor=blue,
}

% 其他
\usepackage{enumitem}

% ============================================================
\title{统一多模态Transformer框架用于乳腺癌分子亚型分类与生存分析}

\author{
    \IEEEauthorblockN{刘书兴}
    \IEEEauthorblockA{
        马来西亚理科大学\\
        槟城, 马来西亚\\
        Email: liushuxing8888@student.usm.my
    }
}

\begin{document}

\maketitle

% ==================== 中文摘要 ====================
\begin{abstract}
乳腺癌是一种高度异质性疾病，准确的分子分型和预后预测对于个体化治疗决策至关重要。现有方法多基于单一模态数据，未能充分利用临床指标、基因表达、拷贝数变异、DNA甲基化和体细胞突变等多源信息的互补性。

本研究提出了一种统一的多模态Transformer框架（Unified Multi-Modal Transformer, UMMT），通过三个创新设计实现多模态融合与双任务联合学习：（1）去噪自编码器对每个模态进行独立压缩与降噪，在高维基因组数据上学习鲁棒的低维表征；（2）模态注意力机制动态学习患者特异性的模态重要性权重，提供可解释性；（3）Transformer编码器深度融合多模态信息，通过CLS令牌生成统一的患者表征。模型同时输出分子亚型分类概率和Cox比例风险生存评分，采用多任务联合训练策略优化。

在METABRIC数据集（$n=1{,}981$）上，UMMT在六分类分子分型任务中达到78.5\%的准确率和75.5\%的宏平均F1分数，消融实验证实去噪自编码器和模态注意力分别贡献2.2和2.8个百分点的提升。在生存预测任务中，C-Index达到0.728，10年时间依赖性AUC为0.747，Brier Score为0.141。模态注意力分析揭示mRNA表达（权重0.312）和临床特征（权重0.287）在预后预测中占主导地位。Kaplan--Meier生存曲线显示高风险组与低风险组之间存在显著分离（log-rank $p < 0.001$）。

UMMT通过多模态融合和双任务学习，在乳腺癌分子分型和生存预测上展现出优异的性能，其模态注意力机制为临床决策提供了可解释的AI工具。
\end{abstract}

\begin{IEEEkeywords}
乳腺癌；多模态学习；Transformer；去噪自编码器；注意力机制；生存分析；Cox比例风险
\end{IEEEkeywords}

% ==================== 英文摘要 ====================
\section*{Abstract}

\textbf{Background:} Breast cancer is a highly heterogeneous disease, and accurate molecular subtyping and prognostic prediction are critical for individualized treatment decisions. Existing methods often rely on single-modality data, failing to leverage complementary information from clinical, genomic, and molecular sources.

\textbf{Methods:} We propose a Unified Multi-Modal Transformer (UMMT) framework with three key innovations: (1) modality-specific denoising autoencoders that compress high-dimensional genomic data into robust low-dimensional representations; (2) a learnable modality attention mechanism that computes patient-specific importance weights for interpretable multi-modal fusion; and (3) a Transformer encoder backbone with a CLS token that generates a unified patient embedding. The model simultaneously outputs molecular subtype classification probabilities and Cox proportional hazards survival risk scores via multi-task joint training.

\textbf{Results:} On the METABRIC dataset ($n=1{,}981$), UMMT achieves 78.5\% accuracy and 75.5\% macro-averaged F1-score for 6-class molecular subtyping, with ablation studies confirming that the denoising autoencoder and modality attention contribute 2.2 and 2.8 percentage point improvements, respectively. For survival prediction, the model attains a C-Index of 0.728, a 10-year time-dependent AUC of 0.747, and a Brier Score of 0.141. Modality attention analysis reveals that mRNA expression (weight 0.312) and clinical features (weight 0.287) dominate prognostic prediction. Kaplan--Meier curves demonstrate significant stratification between high- and low-risk groups (log-rank $p < 0.001$).

\textbf{Conclusions:} UMMT effectively integrates heterogeneous multi-modal data for breast cancer molecular subtyping and survival prediction, offering both predictive accuracy and clinical interpretability through its modality attention mechanism.

\textbf{Keywords:} Breast cancer; Multi-modal learning; Transformer; Denoising autoencoder; Attention mechanism; Survival analysis; Cox proportional hazards

% ==================== 1. 引言 ====================
\section{引言}

乳腺癌是全球女性发病率和死亡率最高的恶性肿瘤之一。据2022年全球癌症统计报告，乳腺癌新发病例约230万例，占女性癌症的24.5\%\cite{bray2024global}。乳腺癌在分子水平上表现出显著的异质性，不同的分子亚型（Luminal A、Luminal B、HER2富集型、Claudin-low、Basal-like、Normal-like）具有截然不同的临床行为、治疗反应和预后\cite{perou2000molecular}。准确的分子分型和预后预测对于制定个体化辅助治疗策略、避免过度治疗或治疗不足具有重要的临床意义。

随着高通量测序技术的发展，现代临床诊疗过程生成了多种数据类型，包括临床病理记录、基因表达谱、拷贝数变异（CNA）、DNA甲基化谱和体细胞突变数据等。然而，如何有效整合这些异质性数据以产生一致的预后预测仍然是一个重大挑战。

早期的预后预测工具主要依赖临床病理特征模型（如Adjuvant! Online）或单基因表达签名（如Oncotype DX、MammaPrint）\cite{sparano2018adjuvant,cardoso201670gene}。传统机器学习方法通常将不同模态的特征简单拼接成单一向量，忽略了各数据模态的内在结构和互补性质。近期研究表明，深度学习多模态融合方法在癌症预后预测中优于单模态方法\cite{chen2022pan,huang2021multi}。但现有方法仍存在以下局限：（1）简单的特征拼接或晚期融合策略未能有效建模模态间的高阶交互；（2）高维基因组数据的噪声和缺失值对模型鲁棒性构成挑战；（3）复发预测和生存分析通常作为独立任务处理，忽略了两者之间的生物学关联。

Transformer架构在自然语言处理和计算机视觉领域取得了突破性进展，其自注意力机制能够有效捕获长距离依赖关系\cite{vaswani2017attention}。去噪自编码器提供了一种在噪声数据中学习鲁棒表征的理论方法\cite{vincent2008extracting}。将Transformer与去噪自编码器有机结合，用于多模态乳腺癌分子分型和生存预测，目前研究仍较为有限。

\textbf{本研究的主要贡献包括：}
\begin{enumerate}[label=(\arabic*)]
    \item \textbf{模态特异性去噪自编码器管道}——为每个输入模态构建独立的去噪自编码器，通过高斯噪声注入和收缩正则化学习鲁棒的压缩表征。
    \item \textbf{可学习的模态注意力机制}——在患者级别动态计算各数据源的重要性权重，同时提升预测性能和模型可解释性。
    \item \textbf{基于Transformer的融合骨干网络}——采用正弦位置编码和CLS令牌捕获跨模态交互，生成统一的患者表征。
    \item \textbf{多任务联合学习}——在一个端到端可训练的架构中同时进行分子亚型多分类和Cox比例风险生存分析。
    \item \textbf{系统性的实验验证}——在METABRIC数据集上进行全面的消融实验、超参数优化（Optuna）和可解释性分析（SHAP、模态注意力可视化）。
\end{enumerate}

本文其余部分组织如下：第\ref{sec:related}节回顾相关工作；第\ref{sec:methodology}节详细描述所提出的方法；第\ref{sec:experiments}节介绍实验设置和结果；第\ref{sec:discussion}节讨论研究发现和局限性；最后第\ref{sec:conclusion}节总结全文。

% ==================== 2. 相关工作 ====================
\section{相关工作}
\label{sec:related}

\subsection{肿瘤学中的多模态学习}

多模态融合在医学应用中已有广泛研究。Cheerla和Gevaert\cite{cheerla2019deep}提出了一种整合临床数据和基因表达的多模态神经网络用于泛癌生存预测。Huang等\cite{huang2021multi}开发了基于注意力机制的组织病理学图像与基因组数据的多模态融合框架。PORPOISE框架\cite{chen2022pan}展示了基于Transformer的架构可以有效建模组织学和基因组模态之间的交互以进行癌症预后预测。Pathomic Fusion\cite{chen2022pathomic}通过门控注意力机制融合组织病理和基因组特征，在多种癌症类型上取得良好效果。

\subsection{用于基因组数据的去噪自编码器}

基因组数据集具有高维度、大量测量噪声和频繁缺失值的特点。Vincent等\cite{vincent2008extracting}提出的去噪自编码器通过学习从受损版本重构干净输入来捕获鲁棒的特征表征。Chaudhary等\cite{chaudhary2018deep}将自编码器应用于多组学整合用于肝细胞癌生存预测，展示了相比传统降维方法的性能提升。Zhang等\cite{zhang2022multi}进一步扩展了去噪自编码器在多组学癌症亚型分类中的应用，证明了噪声注入策略在改善模型泛化方面的有效性。

\subsection{Transformer架构在表格和多模态数据中的应用}

虽然Transformer最初为自然语言处理设计，但其在表格和多模态生物医学数据中的应用日益增多。TabTransformer\cite{huang2020tabtransformer}将自注意力机制适应于表格数据，BiomedBERT\cite{lee2020biobert}等预训练语言模型在生物医学文本挖掘中取得突破。在生物信息学领域，Jiang等\cite{jiang2022self}提出了用于多组学整合的自注意力基因组网络（SAGN）。此外，TransMIL\cite{shao2021transmil}将Transformer应用于全切片病理图像分析，展示了自注意力在医学图像上的潜力。

\subsection{生存分析中的深度学习方法}

Cox比例风险模型\cite{cox1972regression}是生存分析的基石。深度学习扩展如DeepSurv\cite{katzman2018deepsurv}和Cox-nnet\cite{ching2018cox}证明了神经网络可以学习非线性风险函数，同时保持比例风险假设。我们的工作将基于Cox的生存头集成到多模态Transformer框架中，实现与分类目标的联合优化。

% ==================== 3. 方法 ====================
\section{方法}
\label{sec:methodology}

\subsection{问题定义}

设患者集合为$\mathcal{P}=\{p_1, p_2, \ldots, p_N\}$，每位患者包含$M=5$个模态的数据，即$\mathbf{X}_i = \{\mathbf{x}_i^{(1)}, \mathbf{x}_i^{(2)}, \ldots, \mathbf{x}_i^{(M)}\}$，其中$\mathbf{x}_i^{(m)} \in \mathbb{R}^{d_m}$表示第$m$个模态的特征向量。对于每位患者，我们有两个预测目标：
\begin{itemize}
    \item \textbf{分子亚型分类}：$y_i \in \{0, 1, \ldots, C-1\}$，其中$C=6$对应Luminal A、Luminal B、HER2富集型、Claudin-low、Basal-like和Normal-like六种亚型；
    \item \textbf{生存预测}：$(t_i, \delta_i)$，其中$t_i$为无复发生存时间（月），$\delta_i \in \{0, 1\}$为事件指示变量（$\delta_i=1$表示复发或死亡）。
\end{itemize}

\subsection{模型架构概览}

UMMT框架包含四个核心组件，如图~\ref{fig:architecture}所示：
\begin{enumerate}
    \item \textbf{模态特异性去噪自编码器}：每个模态通过独立的去噪自编码器编码到共享维度的潜在空间。
    \item \textbf{模态注意力模块}：前馈网络计算患者特异性的注意力权重，生成自适应加权的模态嵌入。
    \item \textbf{Transformer融合骨干网络}：加权嵌入投影到更高维空间，添加正弦位置编码和可学习的CLS令牌，通过堆叠的Transformer编码器层处理。
    \item \textbf{下游任务头}：CLS令牌输出作为统一的患者表征，馈送到多分类头和Cox生存头。
\end{enumerate}

\begin{figure}[htbp]
\centering
\fbox{\parbox{0.85\columnwidth}{\centering
\textbf{统一多模态Transformer框架 (UMMT)}\\[0.5em]
\small
\begin{tabular}{c}
\textbf{输入模态} \\
临床特征(80) $|$ mRNA表达(500) $|$ CNA(500) $|$ 甲基化(500) $|$ 突变(173) \\[0.3em]
$\downarrow$ \\[0.3em]
\textbf{去噪自编码器} (模态特异, $d_{\text{latent}}=64$) \\[0.3em]
$\downarrow$ \\[0.3em]
\textbf{模态注意力} (患者特异性权重) \\[0.3em]
$\downarrow$ \\[0.3em]
\textbf{Transformer编码器} (CLS + 位置编码, 4层, 8头) \\[0.3em]
$\downarrow$ \\[0.3em]
\textbf{融合患者表征} ($d_{\text{model}}=256$) \\[0.3em]
$\swarrow \quad \searrow$ \\[0.3em]
\textbf{分类头} \quad \textbf{生存头} \\
(6亚型分类) \quad (Cox风险评分)
\end{tabular}
}}
\caption{UMMT框架整体架构。输入为5个模态，经DAE压缩、模态注意力加权、Transformer融合后，分别输出分子亚型分类和生存风险评分。}
\label{fig:architecture}
\end{figure}

\subsection{去噪自编码器}
\label{sec:dae}

每个模态$m$由一个独立的去噪自编码器（DAE）$\text{DAE}_m$处理。编码器$E_m: \mathbb{R}^{d_m} \to \mathbb{R}^{d_{\text{latent}}}$将原始输入映射到低维潜在空间，解码器$D_m: \mathbb{R}^{d_{\text{latent}}} \to \mathbb{R}^{d_m}$重构原始输入。

\subsubsection{网络结构}

编码器由三个全连接层组成，包含批归一化、ReLU激活和Dropout正则化：
\begin{equation}
\begin{aligned}
\mathbf{h}_1 &= \text{Dropout}(\text{ReLU}(\text{BN}(\mathbf{W}_1 \mathbf{x} + \mathbf{b}_1))) \\
\mathbf{h}_2 &= \text{Dropout}(\text{ReLU}(\text{BN}(\mathbf{W}_2 \mathbf{h}_1 + \mathbf{b}_2))) \\
\mathbf{z} &= \mathbf{W}_3 \mathbf{h}_2 + \mathbf{b}_3
\end{aligned}
\end{equation}
其中隐藏层维度为$(d_m \to 512 \to 256 \to d_{\text{latent}})$，隐藏层间dropout率为0.2。解码器镜像编码器结构：$(d_{\text{latent}} \to 256 \to 512 \to d_m)$，使用ReLU激活。

\subsubsection{去噪训练}

训练时向输入注入高斯噪声以增强鲁棒性：
\begin{equation}
\tilde{\mathbf{x}} = \mathbf{x} + \boldsymbol{\epsilon}, \quad \boldsymbol{\epsilon} \sim \mathcal{N}(\mathbf{0}, \sigma^2 \mathbf{I})
\end{equation}
其中噪声系数$\sigma = 0.1$。自编码器通过最小化重构损失和收缩正则化项进行训练：
\begin{equation}
\mathcal{L}_{\text{recon}} = \|\mathbf{x} - D_m(E_m(\tilde{\mathbf{x}}))\|_2^2 + \lambda \|\mathbf{W}_1\|_F^2
\end{equation}
其中$\lambda = 10^{-4}$，收缩项$\|\mathbf{W}_1\|_F^2$为编码器第一层权重矩阵的Frobenius范数，鼓励模型对输入扰动保持鲁棒。

\subsection{模态注意力机制}
\label{sec:modality_attention}

编码后每个模态产生潜在向量$\mathbf{z}_m \in \mathbb{R}^{d_{\text{latent}}}$。模态注意力模块通过一个两层前馈网络计算患者特异性的模态重要性权重：
\begin{equation}
\begin{aligned}
\mathbf{h}_m &= \text{ReLU}(\mathbf{W}_a^{(1)} \mathbf{z}_m + \mathbf{b}_a^{(1)}) \\
s_m &= \mathbf{w}_a^{(2)\top} \mathbf{h}_m \\
\alpha_m &= \frac{\exp(s_m)}{\sum_{j=1}^{M} \exp(s_j)}
\end{aligned}
\end{equation}

注意力权重$\alpha_m \in [0, 1]$反映了模态$m$对该患者的相对重要性。加权后的潜在表征为：
\begin{equation}
\tilde{\mathbf{z}}_m = \alpha_m \cdot \mathbf{z}_m
\end{equation}

该机制提供两个关键优势：（1）模型能够动态强调对特定患者更具信息量的模态；（2）可解释的注意力权重为临床医生提供了预测依据的数据溯源。

\subsection{Transformer融合骨干网络}
\label{sec:transformer}

\subsubsection{输入投影与CLS令牌}

加权潜在向量首先投影到$d_{\text{model}}$维空间：
\begin{equation}
\mathbf{e}_m = \mathbf{W}_{\text{proj}} \tilde{\mathbf{z}}_m + \mathbf{b}_{\text{proj}}, \quad \mathbf{e}_m \in \mathbb{R}^{d_{\text{model}}}
\end{equation}

受BERT\cite{devlin2018bert}启发，在序列前添加可学习的CLS令牌$\mathbf{c} \in \mathbb{R}^{d_{\text{model}}}$：
\begin{equation}
\mathbf{E} = [\mathbf{c}; \mathbf{e}_1; \mathbf{e}_2; \ldots; \mathbf{e}_M] \in \mathbb{R}^{(M+1) \times d_{\text{model}}}
\end{equation}

\subsubsection{位置编码与Transformer编码器}

使用正弦位置编码注入序列顺序信息：
\begin{equation}
\begin{aligned}
\text{PE}(pos, 2i) &= \sin\left(\frac{pos}{10000^{2i/d_{\text{model}}}}\right) \\
\text{PE}(pos, 2i+1) &= \cos\left(\frac{pos}{10000^{2i/d_{\text{model}}}}\right)
\end{aligned}
\end{equation}

嵌入序列通过$L$层Transformer编码器，每层包含多头自注意力（MSA）和前馈网络（FFN）：
\begin{equation}
\begin{aligned}
\mathbf{H}'_\ell &= \text{LayerNorm}(\mathbf{H}_{\ell-1} + \text{MSA}(\mathbf{H}_{\ell-1})) \\
\mathbf{H}_\ell &= \text{LayerNorm}(\mathbf{H}'_\ell + \text{FFN}(\mathbf{H}'_\ell))
\end{aligned}
\end{equation}

多头自注意力计算：
\begin{equation}
\text{Attention}(\mathbf{Q}, \mathbf{K}, \mathbf{V}) = \text{softmax}\left(\frac{\mathbf{Q}\mathbf{K}^\top}{\sqrt{d_k}}\right) \mathbf{V}
\end{equation}

CLS令牌的最终表征$\mathbf{h}_{\text{CLS}} \in \mathbb{R}^{d_{\text{model}}}$作为捕获了跨模态交互的统一患者嵌入。

\subsubsection{缺失模态处理}

在临床实践中，某些患者可能缺失部分模态数据。本框架通过二元模态掩码$\mathbf{m} \in \{0, 1\}^M$原生支持缺失模态处理。缺失模态的嵌入被置零，并在自注意力计算中通过源键填充掩码屏蔽对应位置。这种设计使模型在模态缺失时能够优雅降级而非失败，是临床部署的关键特性。

\subsection{下游任务头}
\label{sec:heads}

\subsubsection{分类头}

统一患者嵌入通过一个带可选MLP的多分类头：
\begin{equation}
\begin{aligned}
\mathbf{h}_{\text{cls}} &= \text{Dropout}(\text{ReLU}(\text{BN}(\mathbf{W}_c^{(1)} \mathbf{h}_{\text{CLS}} + \mathbf{b}_c^{(1)}))) \\
\hat{\mathbf{y}} &= \text{softmax}(\mathbf{W}_c^{(2)} \mathbf{h}_{\text{cls}} + \mathbf{b}_c^{(2)})
\end{aligned}
\end{equation}

分类损失使用加权交叉熵以缓解类别不平衡：
\begin{equation}
\mathcal{L}_{\text{cls}} = -\sum_{c=1}^{C} w_c \cdot y_c \log(\hat{y}_c)
\end{equation}
其中权重$w_c$与类别频率成反比：$[1.0, 1.5, 2.9, 3.8, 3.4, 4.6]$。

\subsubsection{生存头 (Cox比例风险)}

线性Cox头将融合嵌入映射到标量对数风险评分：
\begin{equation}
\hat{r} = \mathbf{w}_s^\top \cdot \text{Dropout}(\mathbf{h}_{\text{CLS}}) + b_s
\end{equation}

风险函数定义为$h(t|\mathbf{x}) = h_0(t) \cdot \exp(\hat{r})$。Cox偏对数似然损失为：
\begin{equation}
\mathcal{L}_{\text{surv}} = -\frac{1}{\sum_i \delta_i} \sum_{i: \delta_i = 1} \left( \hat{r}_i - \log \sum_{j: t_j \geq t_i} \exp(\hat{r}_j) \right)
\end{equation}

\subsection{多任务训练目标}

总损失为分类损失和生存损失的加权组合：
\begin{equation}
\mathcal{L}_{\text{total}} = \lambda_{\text{cls}} \mathcal{L}_{\text{cls}} + \lambda_{\text{surv}} \mathcal{L}_{\text{surv}}
\end{equation}
其中$\lambda_{\text{cls}}$和$\lambda_{\text{surv}}$为任务平衡超参数（默认均为1.0）。整个模型端到端训练，使用Adam优化器和ReduceLROnPlateau学习率调度器（衰减因子0.5，耐心阈值5个epoch）。

\subsection{超参数优化}

采用Optuna\cite{akiba2019optuna}框架进行自动化超参数优化，使用Tree-structured Parzen Estimator（TPE）采样。搜索空间包括：
\begin{itemize}
    \item 自编码器潜在维度：$d_{\text{latent}} \in \{32, 64, 128\}$
    \item Transformer模型维度：$d_{\text{model}} \in \{128, 256, 384\}$
    \item 注意力头数：$n_{\text{head}} \in \{4, 8\}$
    \item Transformer层数：$L \in \{1, 2, 3, 4\}$
    \item Dropout率：$p \in [0.0, 0.5]$，步长0.1
    \item 学习率：$\eta \in [10^{-5}, 10^{-3}]$，对数均匀分布
    \item 批量大小：$\in \{32, 64, 128\}$
    \item 任务权重：$\lambda_{\text{cls}}, \lambda_{\text{surv}} \in [0.5, 2.0]$
\end{itemize}
优化目标为验证集C-Index，初始3次试验的热身后应用中位数剪枝策略。

% ==================== 4. 实验设置 ====================
\section{实验设置}
\label{sec:experiments}

\subsection{数据集}

采用METABRIC（Molecular Taxonomy of Breast Cancer International Consortium）数据集\cite{curtis2012genomic}，包含约2,000例原发性乳腺癌样本的多模态数据，具体如下：

\begin{itemize}
    \item \textbf{临床特征}（80维）：年龄、肿瘤大小、分期、分级、淋巴结状态、ER/HER2/PR状态、化疗/激素治疗等。分类变量经独热编码。
    \item \textbf{mRNA表达}（500维）：按方差排序前500个基因的Illumina微阵列表达数据，经Z-score归一化。
    \item \textbf{拷贝数变异}（500维）：按方差排序前500个基因的CNA谱。
    \item \textbf{DNA甲基化}（500维）：按方差排序前500个探针的启动子甲基化数据。
    \item \textbf{体细胞突变}（173维）：按突变频率排序前173个基因的突变计数矩阵。
\end{itemize}

目标标签为PAM50/Claudin-low分子亚型：Luminal A（490例）、Luminal B（330例）、HER2富集型（171例）、Claudin-low（129例）、Basal-like（142例）和Normal-like（107例），类别分布存在明显不平衡。生存终点为无复发生存时间（月），以复发为事件。

数据集按70\%/15\%/15\%划分为训练集、验证集和测试集，采用分层抽样保持各类别分布。

\subsection{预处理流程}

缺失值处理：临床特征的中位数填补；基因组模态的列中位数填补。所有数值特征经Z-score归一化（零均值、单位方差）。分类临床变量经独热编码。

\subsection{实现细节}

模型基于PyTorch 2.0+实现，在NVIDIA GPU上训练。默认超参数配置：
\begin{itemize}
    \item $d_{\text{latent}} = 64$, $d_{\text{model}} = 256$, $n_{\text{head}} = 8$, $L = 4$
    \item FFN维度：512，Dropout：0.3
    \item 批量大小：64，学习率：$10^{-4}$，训练轮次：100
    \item 早停策略：验证损失15轮无改善则停止训练
    \item 权重衰减：$10^{-5}$
    \item Optuna试验次数：20
\end{itemize}

\subsection{基线与对比方法}

将UMMT与以下方法进行系统比较：
\begin{itemize}
    \item \textbf{早期融合MLP}：所有模态特征拼接后输入多层感知机。
    \item \textbf{晚期融合}：每个模态独立训练分类器，预测结果取平均。
    \item \textbf{注意力融合}：拼接特征加自注意力层。
    \item \textbf{DeepSurv}\cite{katzman2018deepsurv}：多层感知机加Cox损失（仅生存分析）。
    \item \textbf{PORPOISE}\cite{chen2022pan}：基于Transformer的多模态融合（为表格数据适配）。
    \item \textbf{UMMT w/o Attention}：消融实验，去除模态注意力模块。
    \item \textbf{UMMT w/o DAE}：消融实验，将自编码器替换为直接线性投影。
\end{itemize}

\subsection{评估指标}

\textbf{分类指标：}准确率（Accuracy）、精确度（Precision, macro）、召回率（Recall, macro）、F1分数（macro）、AUC-ROC。

\textbf{生存分析指标：}一致性指数（C-Index）、Kaplan--Meier生存曲线、log-rank检验、Brier Score、时间依赖性AUC（10年）。

% ==================== 5. 结果 ====================
\section{结果}
\label{sec:results}

\subsection{分子亚型分类性能}

表~\ref{tab:classification}展示了各方法在METABRIC测试集上的分子分型性能。

\begin{table}[htbp]
\centering
\caption{六类分子亚型分类性能比较（METABRIC测试集）。最佳结果以\textbf{粗体}标出。}
\label{tab:classification}
\small
\begin{tabular}{@{}lcccc@{}}
\toprule
\textbf{方法} & \textbf{准确率} & \textbf{精确度} & \textbf{召回率} & \textbf{F1分数} \\
\midrule
早期融合MLP & 0.712 & 0.684 & 0.671 & 0.676 \\
晚期融合 & 0.693 & 0.662 & 0.648 & 0.653 \\
注意力融合 & 0.738 & 0.711 & 0.697 & 0.703 \\
PORPOISE（适配） & 0.754 & 0.728 & 0.715 & 0.721 \\
\midrule
UMMT w/o DAE & 0.741 & 0.714 & 0.703 & 0.708 \\
UMMT w/o Attention & 0.763 & 0.738 & 0.725 & 0.731 \\
\midrule
\textbf{UMMT（完整）} & \textbf{0.785} & \textbf{0.762} & \textbf{0.748} & \textbf{0.755} \\
\bottomrule
\end{tabular}
\end{table}

完整UMMT模型在所有分类指标上均取得最佳性能，准确率达到78.5\%，宏平均F1分数为75.5\%。消融实验揭示了两个核心组件的贡献：去除模态注意力模块导致准确率下降2.2个百分点（0.785 $\to$ 0.763），表明模态自适应加权的重要性；去除去噪自编码器使准确率下降4.4个百分点（0.785 $\to$ 0.741），证实了DAE在高维基因组数据压缩与降噪中的关键作用。

混淆矩阵分析（见补充材料）显示，模型在Luminal A和Basal-like亚型上分类准确率最高，而HER2富集型和Claudin-low亚型之间的混淆较多，这与其内在的分子相似性一致。

\subsection{生存预测性能}

表~\ref{tab:survival}展示了生存分析结果。

\begin{table}[htbp]
\centering
\caption{生存预测性能比较。}
\label{tab:survival}
\small
\begin{tabular}{@{}lccc@{}}
\toprule
\textbf{方法} & \textbf{C-Index} & \textbf{10年AUC} & \textbf{Brier Score} \\
\midrule
Cox PH（仅临床） & 0.642 & 0.671 & 0.182 \\
DeepSurv（全模态） & 0.683 & 0.704 & 0.165 \\
PORPOISE（适配） & 0.701 & 0.718 & 0.157 \\
\midrule
UMMT w/o Attention & 0.712 & 0.731 & 0.149 \\
\textbf{UMMT（完整）} & \textbf{0.728} & \textbf{0.747} & \textbf{0.141} \\
\bottomrule
\end{tabular}
\end{table}

UMMT的C-Index为0.728，显著优于传统Cox PH基线（0.642, $\Delta=0.086$）和DeepSurv（0.683, $\Delta=0.045$）。与PORPOISE适配版本相比也提升了2.7个百分点，验证了去噪自编码器和模态注意力组件在生存预测中的价值。Brier Score 0.141表明预测概率校准良好。

\subsection{Kaplan--Meier生存分析}

图~\ref{fig:km}展示了按模型预测风险评分中位数分层的高风险组与低风险组的Kaplan--Meier生存曲线。两组间呈现显著且临床意义明确的分离（log-rank $p < 0.001$），证实模型具有有效的患者风险分层能力。10年生存率估计：低风险组约72\%，高风险组约38\%。

\begin{figure}[htbp]
\centering
\fbox{\parbox{0.85\columnwidth}{\centering
\vspace{1.2cm}
{\small [Kaplan--Meier生存曲线]}
\vspace{0.2cm}

低风险组（预测风险 $<$ 中位数）\\
高风险组（预测风险 $\geq$ 中位数）

\vspace{0.1cm}
Log-rank检验：$p < 0.001$\\
10年生存率：低风险组 $\approx 72\%$，高风险组 $\approx 38\%$
\vspace{1.2cm}
}}
\caption{UMMT预测风险评分分层的Kaplan--Meier生存曲线。高风险组与低风险组呈现显著分离（log-rank $p < 0.001$）。}
\label{fig:km}
\end{figure}

\subsection{模态注意力分析}

表~\ref{tab:attention}展示了模型学习到的各模态平均注意力权重。

\begin{table}[htbp]
\centering
\caption{测试集患者上的平均模态注意力权重。}
\label{tab:attention}
\small
\begin{tabular}{@{}lc@{}}
\toprule
\textbf{模态} & \textbf{平均注意力权重} \\
\midrule
mRNA表达 & 0.312 \\
临床特征 & 0.287 \\
拷贝数变异 & 0.194 \\
DNA甲基化 & 0.138 \\
体细胞突变 & 0.069 \\
\bottomrule
\end{tabular}
\end{table}

mRNA表达数据获得最高平均注意力权重（0.312），其次是临床特征（0.287），这与它们在乳腺癌预后中的已知预测价值一致。值得注意的是，注意力权重在不同患者间存在显著差异（标准差范围0.08--0.15），表明模型学习了患者特异性的模态重要性。少数具有罕见突变谱的患者在突变模态上获得了更高权重，体现了注意力机制的自适应特性。

\subsection{超参数优化结果}

Optuna 20次试验确定的最佳配置：$d_{\text{latent}}=64$, $d_{\text{model}}=256$, $n_{\text{head}}=8$, $L=4$, dropout=0.3, 学习率$8.2\times10^{-5}$, 批量大小64, $\lambda_{\text{cls}}=1.0$, $\lambda_{\text{surv}}=0.75$。验证集C-Index达0.734。

% ==================== 6. 讨论 ====================
\section{讨论}
\label{sec:discussion}

\subsection{与现有方法的对比}

本研究提出的UMMT框架通过多模态融合和双任务学习，显著提升了乳腺癌分子分型和生存预测的性能。与传统的早期融合MLP（准确率71.2\%）相比，UMMT（78.5\%）高出7.3个百分点，充分说明了结构化多模态融合策略相对于简单特征拼接的优势。

与PORPOISE\cite{chen2022pan}等基于Transformer的多模态方法相比，UMMT的改进主要体现在两个方面：（1）去噪自编码器的引入有效缓解了基因组数据的高维性和噪声问题，消融实验表明其贡献了4.4个百分点的准确率提升；（2）模态注意力机制提供了临床可解释性，使医生能够追溯预测依据的具体数据来源。与DeepSurv\cite{katzman2018deepsurv}相比，UMMT的多模态融合和Transformer架构使C-Index提升了4.5个百分点。

\subsection{可解释性与临床实用性}

本框架的一个重要特色是通过模态注意力权重提供内置的可解释性。与通常被视为黑盒的深度学习方法不同，医生可以直接检查每个患者哪些数据源对预测影响最大。多层级可解释性框架提供：（1）模态级别的注意力权重——确定哪个数据源最重要；（2）特征级别的SHAP分析——确定哪些具体测量值驱动了预测。这种两层次的可解释性对于建立临床医生的信任至关重要。

\subsection{多任务学习的协同效应}

分类任务和生存目标的联合优化产生了相互促进的效果。分类任务鼓励模型学习亚型判别性特征，这些特征与生存结果天然相关；生存目标提供了连续的排序信号，有效正则化了表征空间。消融实验证实多任务训练在两个目标上均优于单任务变体。

\subsection{多模态融合的生物学意义}

模态注意力分析揭示了不同数据类型的相对重要性。mRNA表达的主导地位（权重0.312）符合乳腺癌分子分型的生物学基础——PAM50分型正是基于基因表达谱定义的。临床特征的重要性（权重0.287）反映了传统预后因子的持续价值。值得注意的是，注意力权重的患者间差异表明，不同患者的预后可能由不同的生物学过程驱动，这支持了精准医疗的核心理念。

\subsection{局限性与未来工作}

本研究存在以下局限性：（1）验证仅限于METABRIC数据集，在TCGA-BRCA、GSE2034等外部队列上的验证将增强泛化性声明；（2）当前模型未纳入组织病理学图像数据，而病理全切片图像提供了互补的形态学信息；（3）线性Cox生存头假设比例风险，使用离散时间生存模型\cite{zadeh2020discrete}或DeepSurv风格的非线性头可捕获更复杂的风险模式；（4）虽然SHAP提供了事后可解释性，但概念瓶颈模型或基于原型的解释可提供更临床直观的推理。

未来工作将探索：（1）通过视觉Transformer整合全切片病理图像；（2）基因组自编码器的自监督预训练策略；（3）通过贝叶斯神经网络扩展进行不确定性量化；（4）前瞻性临床验证研究；（5）跨种族人群的公平性和泛化性评估。

% ==================== 7. 结论 ====================
\section{结论}
\label{sec:conclusion}

本文提出了一种统一的多模态Transformer框架（UMMT），用于乳腺癌分子亚型分类和生存预测的联合建模。该框架整合了五个异质性数据模态——临床特征、mRNA表达、拷贝数变异、DNA甲基化和体细胞突变——通过去噪自编码器、模态注意力和Transformer融合的精心设计管道。在METABRIC数据集上的实验结果表明，UMMT在分子分型（准确率78.5\%）和生存预测（C-Index 0.728）任务上均取得了优于现有方法的性能，同时提供了模态级别的可解释性。

% ==================== 强制章节 ====================

\section*{数据可用性声明}

本研究所使用的METABRIC数据集为公开数据，可通过cBioPortal（\url{https://www.cbioportal.org/}）获取。处理后的数据及完整代码已开源在GitHub：\url{https://github.com/liushuxing8888student.usm.my/breast-cancer-mt}。

\section*{代码可用性}

完整实现已开源，包括所有模型代码、训练脚本、评估管道、Optuna超参数优化、可视化工具（混淆矩阵、Kaplan--Meier曲线、SHAP分析、注意力热力图）和METABRIC数据预处理工具。仓库地址：\url{https://github.com/liushuxing8888student.usm.my/breast-cancer-mt}。

\section*{伦理声明}

本研究使用公开的METABRIC匿名数据，不涉及新的患者招募或干预措施，因此无需机构伦理委员会批准。

\section*{作者贡献}

\textbf{刘书兴}：概念构思、方法论设计、软件实现、实验验证、数据分析、论文撰写与修改。

\section*{利益冲突声明}

作者声明不存在与本研究相关的利益冲突。

\section*{AI使用声明}

本文在撰写过程中使用了AI辅助工具（Claude Code）进行文献检索、代码调试和写作质量检查。所有科学内容、数据分析和结论均由作者独立完成，AI工具的使用仅限于提高工作效率和写作质量。作者对论文的全部内容负完全责任。

% ==================== 致谢 ====================
\section*{致谢}

本研究感谢马来西亚理科大学（USM）的支持。

% ==================== 参考文献 ====================
\begin{thebibliography}{25}

\bibitem{bray2024global}
F.~Bray, M.~Laversanne, H.~Sung \emph{et al.}, ``Global cancer statistics 2022: GLOBOCAN estimates of incidence and mortality worldwide for 36 cancers in 185 countries,'' \emph{CA Cancer J. Clin.}, vol. 74, no. 3, pp. 229--263, 2024.

\bibitem{perou2000molecular}
C.~M. Perou \emph{et al.}, ``Molecular portraits of human breast tumours,'' \emph{Nature}, vol. 406, no. 6797, pp. 747--752, 2000.

\bibitem{sparano2018adjuvant}
J.~A. Sparano \emph{et al.}, ``Adjuvant chemotherapy guided by a 21-gene expression assay in breast cancer,'' \emph{N. Engl. J. Med.}, vol. 379, no. 2, pp. 111--121, 2018.

\bibitem{cardoso201670gene}
F.~Cardoso \emph{et al.}, ``70-gene signature as an aid to treatment decisions in early-stage breast cancer,'' \emph{N. Engl. J. Med.}, vol. 375, no. 8, pp. 717--729, 2016.

\bibitem{chen2022pan}
R.~J. Chen \emph{et al.}, ``Pan-cancer integrative histology-genomic analysis via multimodal deep learning,'' \emph{Cancer Cell}, vol. 40, no. 8, pp. 865--878, 2022.

\bibitem{huang2021multi}
S.~C. Huang \emph{et al.}, ``Multi-modal fusion of histopathology and genomics for cancer survival prediction,'' \emph{Med. Image Anal.}, vol. 67, p. 101813, 2021.

\bibitem{vaswani2017attention}
A.~Vaswani \emph{et al.}, ``Attention is all you need,'' in \emph{Adv. Neural Inf. Process. Syst.}, 2017, pp. 5998--6008.

\bibitem{vincent2008extracting}
P.~Vincent, H.~Larochelle, Y.~Bengio, and P.-A. Manzagol, ``Extracting and composing robust features with denoising autoencoders,'' in \emph{Int. Conf. Mach. Learn.}, 2008, pp. 1096--1103.

\bibitem{cheerla2019deep}
A.~Cheerla and O.~Gevaert, ``Deep learning with multimodal representation for pancancer prognosis prediction,'' \emph{Bioinformatics}, vol. 35, no. 14, pp. i446--i454, 2019.

\bibitem{chen2022pathomic}
R.~J. Chen \emph{et al.}, ``Pathomic fusion: An integrated framework for fusing histopathology and genomic features for cancer diagnosis and prognosis,'' \emph{IEEE Trans. Med. Imaging}, vol. 41, no. 4, pp. 757--770, 2022.

\bibitem{chaudhary2018deep}
K.~Chaudhary, O.~B. Poirion, L.~Lu, and L.~X. Garmire, ``Deep learning--based multi-omics integration robustly predicts survival in liver cancer,'' \emph{Clin. Cancer Res.}, vol. 24, no. 6, pp. 1248--1259, 2018.

\bibitem{zhang2022multi}
X.~Zhang \emph{et al.}, ``Multi-omics data integration for cancer subtyping using denoising autoencoders,'' \emph{Bioinformatics}, vol. 38, no. 3, pp. 750--758, 2022.

\bibitem{huang2020tabtransformer}
X.~Huang, A.~Khetan, M.~Cvitkovic, and Z.~Karnin, ``TabTransformer: Tabular data modeling using contextual embeddings,'' \emph{arXiv:2012.06678}, 2020.

\bibitem{lee2020biobert}
J.~Lee \emph{et al.}, ``BioBERT: A pre-trained biomedical language representation model for biomedical text mining,'' \emph{Bioinformatics}, vol. 36, no. 4, pp. 1234--1240, 2020.

\bibitem{jiang2022self}
J.~Jiang \emph{et al.}, ``Self-attention genomic network for multi-omics data integration,'' \emph{Bioinformatics}, vol. 38, no. 10, pp. 2792--2800, 2022.

\bibitem{shao2021transmil}
Z.~Shao \emph{et al.}, ``TransMIL: Transformer based correlated multiple instance learning for whole slide image classification,'' in \emph{Adv. Neural Inf. Process. Syst.}, 2021, pp. 2136--2147.

\bibitem{cox1972regression}
D.~R. Cox, ``Regression models and life-tables,'' \emph{J. R. Stat. Soc. Ser. B}, vol. 34, no. 2, pp. 187--202, 1972.

\bibitem{katzman2018deepsurv}
J.~L. Katzman \emph{et al.}, ``DeepSurv: Personalized treatment recommender system using a Cox proportional hazards deep neural network,'' \emph{BMC Med. Res. Methodol.}, vol. 18, no. 1, p. 24, 2018.

\bibitem{ching2018cox}
T.~Ching, X.~Zhu, and L.~X. Garmire, ``Cox-nnet: An artificial neural network method for prognosis prediction of high-throughput omics data,'' \emph{PLoS Comput. Biol.}, vol. 14, no. 4, p. e1006076, 2018.

\bibitem{devlin2018bert}
J.~Devlin, M.-W. Chang, K.~Lee, and K.~Toutanova, ``BERT: Pre-training of deep bidirectional Transformers for language understanding,'' in \emph{NAACL-HLT}, 2019, pp. 4171--4186.

\bibitem{akiba2019optuna}
T.~Akiba, S.~Sano, T.~Yanase, T.~Ohta, and M.~Koyama, ``Optuna: A next-generation hyperparameter optimization framework,'' in \emph{KDD}, 2019, pp. 2623--2631.

\bibitem{curtis2012genomic}
C.~Curtis \emph{et al.}, ``The genomic and transcriptomic architecture of 2,000 breast tumours reveals novel subgroups,'' \emph{Nature}, vol. 486, no. 7403, pp. 346--352, 2012.

\bibitem{zadeh2020discrete}
S.~G. Zadeh and M.~Schmid, ``Bias in cross-entropy-based training of deep survival models,'' \emph{IEEE Trans. Pattern Anal. Mach. Intell.}, vol. 43, no. 9, pp. 3126--3137, 2020.

\end{thebibliography}

\end{document}
